\input{../tex-inputs/latex-leseansicht-vorspann}

\section[Arthur Schnitzler, Richard Beer-Hofmann und Paul Goldmann an Gustav Schwarzkopf, 23. 8. 1900]{L04506 Arthur Schnitzler, Richard Beer-Hofmann und Paul Goldmann an Gustav
               Schwarzkopf, 23. 8. 1900}
\nopagebreak\mylabel{L04506v}
\rehead{ }\normalsize\beginnumbering\briefempfaengerindex{Schwarzkopf, Gustav@\textsc{Schwarzkopf, Gustav}!zzzGoldmann, Paul@\emph{von Paul Goldmann}!1900-08-231@{23. 8. 1900}|(be}\briefempfaengerindex{Schwarzkopf, Gustav@\textsc{Schwarzkopf, Gustav}!zzzBeer-Hofmann, Richard@\emph{von Richard Beer-Hofmann}!1900-08-231@{23. 8. 1900}|(be}\briefempfaengerindex{Schwarzkopf, Gustav@\textsc{Schwarzkopf, Gustav}!zzzSchnitzler, Arthur@\emph{von Arthur Schnitzler}!1900-08-231@{23. 8. 1900}|(be}
\toendnotes[C]{\smallbreak\pagebreak[2]}
\correspDesc{Versand  durch Arthur Schnitzler, Richard Beer-Hofmann, Paul Goldmann am 23. 8. 1900 in Pontresina
\newline{}Übermittlung  am 23. 8. 1900 in Pontresina
\newline{}Zustellung  am 25. 8. 1900 in Wien
\newline{}Erhalt  durch Gustav Schwarzkopf im Zeitraum [25. 8. 1900
                  – 28. 8. 1900?] in Wien}\toendnotes[C]{\smallbreak}
\Standort{DLA, A:Schnitzler, HS.1985.1.1897.}
\physDesc{Bildpostkarte, 136 Zeichen
\newline{}Handschrift Arthur Schnitzler: Bleistift, deutsche Kurrent
\newline{}Handschrift Richard Beer-Hofmann: Bleistift, lateinische Kurrent
\newline{}Handschrift Paul Goldmann: Bleistift, lateinische Kurrent
\newline{}Versand: 1) Stempel: »\nobreak{}\oindex{Pontresina@\textbf{Pontresina}|pwk}Pontresina, 23. VIII. 00., 7\nobreak{}«.   2) Stempel: »\nobreak{}\oindex{I., Innere Stadt@\textbf{I., Innere Stadt}|pwk}Wien 1/1, 25/\textcolor{gray}{8} 00, 9–\textcolor{gray}{10½} V., bestellt\nobreak{}«. }\toendnotes[C]{\smallbreak}
\pstart
           \noindent{}\centering{}{\pb}\textcolor{gray}{\textbf{PONTRESINA}}\oindex{Pontresina@\textbf{Pontresina}|pw}\pend
           \pstart{}{\pb}Hrn \textsc{Gustav Schwarzkopf}\pend{}\pstart{}Wien\oindex{Wien@\textbf{Wien}|pw}\pend{}\pstart{}\textsc{I. Tiefer Graben 23}\oindex{Wien@\textbf{Wien}!I., Innere Stadt@\textbf{I., Innere Stadt}!Tiefer Graben 23@\textbf{Tiefer Graben 23}|pw}\pend{}{\bigskip}\vspace{1em}
\pstart
           \noindent{}{\pb}\label{T_L04506-1v}\edtext{Mit herzlichen Grüßen!}{\lemma{\textnormal{\emph{Mit herzlichen Grüßen!}}}\Cendnote{\textnormal{Die Beschriftung findet sich rechts neben dem Bild seitlich aufgestellt.}}}\label{T_L04506-1}\pend
           
\pstart
           \spacefill\mbox{ArthSch}{\\[\baselineskip]}{[}hs. Beer-Hofmann:{]} \spacefill\mbox{Richard}{\\[\baselineskip]}{[}hs. Goldmann:{]} \spacefill\mbox{Dr. Paul Goldmann.}\pend
           \leftskip=0em{}
\pstart
           \noindent{}{[}hs. Schnitzler:{]} (Frl. G.\pwindex{Grünwald, Ida 28.\,6.\,1873 Wien – Mai 1908 Alland@\textsc{Grünwald, Ida} (28.\,6.\,1873 Wien – Mai 1908 Alland)|pw} wird
                     \label{K_L04506-1v}\edtext{die Novelle\pwindex{Schnitzler, Arthur 15. 5. 1862 Wien – 21. 10. 1931 ebd.@\textsc{Schnitzler, Arthur} (15. 5. 1862 Wien – 21. 10. 1931 ebd.)!Frau Bertha Garlan. Roman@\strich\emph{Frau Bertha Garlan. Roman}|pwv}}{\lemma{\textnormal{\emph{die Novelle}}}\Cendnote{\textnormal{Im Brief vom XXXX Auszeichnungsfehler: Dokument L04080 nicht gefunden erwähnte Schnitzler, dass er den Text von \emph{Frau
                        Bertha Garlan}\pwindex{Schnitzler, Arthur 15. 5. 1862 Wien – 21. 10. 1931 ebd.@\textsc{Schnitzler, Arthur} (15. 5. 1862 Wien – 21. 10. 1931 ebd.)!Frau Bertha Garlan. Roman@\strich\emph{Frau Bertha Garlan. Roman}|pwk} durch seine Schreibkraft Ida Grünwald\pwindex{Grünwald, Ida 28.\,6.\,1873 Wien – Mai 1908 Alland@\textsc{Grünwald, Ida} (28.\,6.\,1873 Wien – Mai 1908 Alland)|pwk} an Schwarzkopf\pwindex{Schwarzkopf, Gustav 7.\,11.\,1853 Wien – 13.\,11.\,1939 ebd.@\textsc{Schwarzkopf, Gustav} (7.\,11.\,1853 Wien – 13.\,11.\,1939 ebd.)|pwk} hatte überbringen lassen.}}}\label{K_L04506-1} holen.)\pend
           
\pstart
           23. 8. 900\pend
           \selectlanguage{ngerman}\endnumbering\briefempfaengerindex{Schwarzkopf, Gustav@\textsc{Schwarzkopf, Gustav}!zzzGoldmann, Paul@\emph{von Paul Goldmann}!1900-08-231@{23. 8. 1900}|)be}\briefempfaengerindex{Schwarzkopf, Gustav@\textsc{Schwarzkopf, Gustav}!zzzBeer-Hofmann, Richard@\emph{von Richard Beer-Hofmann}!1900-08-231@{23. 8. 1900}|)be}\briefempfaengerindex{Schwarzkopf, Gustav@\textsc{Schwarzkopf, Gustav}!zzzSchnitzler, Arthur@\emph{von Arthur Schnitzler}!1900-08-231@{23. 8. 1900}|)be}\mylabel{L04506h}
\begin{anhang}
\end{anhang}\newcommand{\dateiname}{L04506}\newcommand{\titel}{Arthur Schnitzler, Richard Beer-Hofmann und Paul Goldmann an Gustav Schwarzkopf, 23. 8. 1900}\newcommand{\editorInnen}{Herausgegeben von Jahnke, SelmaMüller, Martin Anton}\input{../tex-inputs/latex-leseansicht-abspann}
